% !TEX root = tarasov-qasp-zh.tex
% 章首扉图说明：第 15 章位于 sec-07 节内，本章（照英文源 ch-15.tex）章首无 sec 扉图块；
% 本章尾部有 sec-08 扉图块（服务于下一节），照英文源位置译于本文件末尾（先例：zh/ch-07-zh.tex 尾部）。
\chapter[怎样确定重心的位置？]{怎样确定重心的位置？\label{ch-15}}

\begin{dialogue}

\tchr 在许多情形里，考生都难以确定单个物体或物体系统的重心。这方面的问题你都完全清楚了吗？
\pnote{重心的概念在中学教材里两次给出：人教版初中物理八年级下册“力”一章的“重力”一节指出，质地均匀、外形规则的物体的重心在它的几何中心；高中必修第一册第三章“重力与弹力”一节再次给出（重力的等效作用点）。但组合图形（下文的 L 形平板、带孔圆盘）与多个重物组成的系统，其重心计算超出中学教材要求，属物理竞赛与大学普通物理（静力学）的常规内容；本章先后用到的正是两类基本技巧——把平板分割成矩形（分割法），以及“补上”一个负重量的小圆来处理圆孔（负面积法，见下）。}

\stda 不能说完全清楚。\figr{fig-63}~\drkgry{(a)} 和\figr{fig-64}~\drkgry{(a)} 所示两种情形中的重心该怎样求，我就不很明白。

\begin{figure}[!htb]%[3cm]
\centering
\includegraphics[width=0.4\textwidth]{figs/ch-15/fig-063a.pdf}
\caption[求所给物体的重心。]{问题是求所给物体的重心。\label{fig-63}}
\end{figure}

\tchr 好。第一种情形，照\figr{fig-63}~\drkgry{(b)} 中的虚线把平板分成两个矩形较为方便。矩形 1 的重心在点 $A$；这块矩形的重量与它的面积成正比，由图显然可见，等于 6 个单位（这里重量约定以平方厘米来计量）。矩形 2 的重心在点 $B$；这块矩形的重量等于 10 个单位。然后把点 $A$ 和 $B$ 投影到坐标轴 $O_{x}$ 和 $O_{y}$ 上；在 $X$ 轴上的投影记作 $A_{1}$ 和 $B_{1}$，在 $Y$ 轴上的投影记作 $A_{2}$ 和 $B_{2}$。接着考察“杆” $A_{1}B_{1}$ 和 $A_{2}B_{2}$，假定质量集中在“杆”的两端，每端的质量等于相应矩形的质量（见\figr{fig-63}~\drkgry{(b)}）。这样，求平板重心的问题就化成了求“杆” $A_{1}B_{1}$ 和 $A_{2}B_{2}$ 的重心；这两个重心的位置便是平板重心的坐标。
\pnote{底本此处有排印脱漏：``The centre of gravity of rectangle 2 is at point the weight of this rectangle is equal to 10 units.''，``at point'' 后显然脱 ``$B$;''（下文正要对 $A$、$B$ 两点取投影），中文按“矩形 2 的重心在点 $B$”补足译出。另：底本本章还有几处拼写与标点级排印小误（首句 ``a body of system of bodies'' 中第二个 ``of'' 应为 ``or''、``sup porting'' 断字、两处括号未闭合、一处句末脱句点等），均径改不注。}

不过，让我们把这道题算完。先用众所周知的力矩定则确定“杆” $A_{1}B_{1}$ 的重心位置（见\figr{fig-63}~\drkgry{(b)}）：
\pnote{这条“力矩定则”就是初中学过的杠杆平衡条件（人教版初中物理八年级下册“杠杆”一节：$F_{1}L_{1}=F_{2}L_{2}$），也就是\chap{ch-14} 中的平衡第二条件（力矩条件）用于两个平行力的情形：合力作用点到两力作用点的距离，与两力的大小成反比。$6x=10(2-x)$ 正是按它列出的。}
\begin{equation*}%
6x= 10(2-x) \; \text{由此} \;\; x=\dfrac{5}{4} \, \si{\centi\meter}
\end{equation*}
于是，在所选的坐标系中，平板重心的 $X$ 坐标为
\begin{equation*}%
X = (1 +x) \, \si{\centi\meter}= \frac{9}{4} \, \si{\centi\meter}
\end{equation*}
用同样的方法求“杆” $A_{2}B_{2}$ 的重心：\\
 \begin{equation*}
 6y=10(1-y)
\end{equation*}
 由此得 $y=\dfrac{5}{8} \; \si{\centi\meter}$。于是平板重心的 $Y$ 坐标为
\begin{equation*}%
Y = (1.5+y) \si{\centi\meter}= \frac{17}{8} \si{\centi\meter}
\end{equation*}

\stda 现在我明白了。求平板重心的坐标 $X$，我正是这样做的；只是我不敢肯定坐标 $Y$ 也能用同样的方法求出。

\tchr 来看\figr{fig-64}~\drkgry{(a)} 所示的第二种情形。
\begin{figure}
\centering
\includegraphics[width=0.6\textwidth,angle=2]{figs/ch-15/fig-064a.pdf}
\caption[求所给物体的重心。]{问题是求所给物体的重心。\label{fig-64}}
\end{figure}
这道题有两种解法。比如，可以不直接处理所给的带一个圆孔的圆，而处理由两个物体组成的系统：一个带两个对称圆孔的圆盘，和一个嵌入其中一孔的小圆（\figr{fig-64}~\drkgry{(b)}）。这两个物体的重心都位于各自的几何中心。已知带两孔圆盘的重量与它的面积成正比，即
\begin{equation*}%
\left ( \pi R^{2} -  \frac{2 \pi R^{2}}{4} \right) = \frac{\pi R^{2}}{2}
\end{equation*}
而小圆的重量与它的面积 $ \pi R^{2}/4 $ 成正比，于是问题化为求\figr{fig-64}~\drkgry{(b)} 下方所示两个平行力的合力作用点。用 $x$ 表示所求重心到大圆几何中心的距离。由\figr{fig-64}~\drkgry{(b)}，可以写出
\begin{equation*}%
\left( \frac{\pi R^{2}}{4} \right) \left( \frac{R}{2} -x \right) = \frac{\pi R^{2}}{2} x \;\; \text{由此} \;\; x=\frac{R}{6}
\end{equation*}
还有另一种可能的解法：把所给的带孔圆盘换成一块实心圆盘（无孔），再加上一个位于圆孔原处、具有负重量（negative weight，即方向朝上）的圆 \drkgry{(\figr{fig-64}~(c))}，由它抵消实心圆盘相应部分的正重量。整体看来，这一组合就相当于最初那个带圆孔的圆盘。
\pnote{第二种解法即竞赛与教辅中惯称的“负面积（负质量）法”，是处理均匀物体挖孔问题的标准技巧；“负重量”并不是真实存在的力，纯属代数手段。两种解法最后都化为对同一个未知点取矩的方程，结果自然相同（$x=R/6$）。本章章末习题正是这类挖孔题。}

这时，问题同样化为求\figr{fig-64}~\drkgry{(c)} 底部所示两个力的合力作用点。由图可以写出：
\begin{equation*}%
\pi R^{2} x = \left( \frac{\pi R^{2}}{4} \right) \left( \frac{R}{2}+x \right)
\end{equation*}
由此，与前一种情形一样，$x=\dfrac{R}{6}$。

\stda 第一种解法我更喜欢些，因为它不需要引入负重量。

\tchr 另外，我想提出\hlred{一道题：求\figr{fig-65}~\drkgry{(a)} 所示重物系统的重心。给定六个重量不同的重物（$P_{1},P_{2}, \ldots{}, P_{6}$），沿一根杆排列，彼此距离相等，均为 $a$。杆的重量忽略不计。这道题你打算怎样解？}
\begin{figure}[!htb]%[-6cm]%
\centering
\includegraphics[width=0.6\textwidth,angle=-2.5]{figs/ch-15/fig-065a.pdf}
\caption[求给定配置中重物系统的重心。]{问题是求给定配置中重物系统的重心。\label{fig-65}}
\end{figure}

\stda 我会先取两个重物，比如 $P_{1}$ 和 $P_{2}$，求出它们合力的作用点；再把这个合力（等于 $P_{1} + P_{2}$）画在图上，并把力 $P_{1}$ 和 $P_{2}$ 划掉，不再考虑。这样，原来的六个力只剩五个。接着再求另一对力的合力作用点，依此类推。如此逐次操作，最终就能求出所需的总合力，它的作用点就是整个系统的重心。

\tchr 你的解法完全正确，但实在太繁琐了。我可以给你看一个优雅得多的解法。我们先假定把系统支承在它的重心上（即\figr{fig-65}~\drkgry{(b)} 中的点 $B$）。

\stdb （插话）可是您还不知道重心的位置。您怎么会知道它在力 $P_{3}$ 和 $P_{4}$ 的作用点之间呢？

\tchr 重心究竟在哪里，对我无关紧要。\figr{fig-65}~\drkgry{(b)} 中重心恰好落在力 $P_{3}$ 和 $P_{4}$ 的作用点之间，这一事实我并不利用。就这样，我们假定系统被支承在它的重心上。结果，杆处于平衡状态。除了那六个力之外，还有一个力——支承力 $N$——作用在杆上。既然杆处于平衡状态，就可以应用平衡条件（见\chap{ch-14}）。先对各个力沿竖直方向的投影写出第一个平衡条件：
\begin{equation}%
N=P_{1}+P_{2}+P_{3}+P_{4}+P_{5}+P_{6}
\label{eq-86}
\end{equation}
然后应用第二个条件（力矩条件），对\figr{fig-65}~\drkgry{(b)} 中的点 $A$（即杆的左端）考虑各力的力矩。这里，各力都力图使杆顺时针转动，而支承力力图使它逆时针转动。可以写出
\begin{equation}%
N (\overline{AB}) = aP_{2} + 2aP_{3} + 3aP_{4} + 4aP_{5} + 5aP_{6}
\label{eq-87}
\end{equation}
把条件~\eqref{eq-86} 与~\eqref{eq-87} 联立，即可求出长度 $\overline{AB}$，即从杆左端量起的所求重心位置：
\begin{equation}%
\overline{AB} = \frac{aP_{2} + 2aP_{3} + 3aP_{4} + 4aP_{5} + 5aP_{6}}{P_{1}+P_{2}+P_{3}+P_{4}+P_{5}+P_{6}}
\label{eq-88}
\end{equation}
\pnote{式~\eqref{eq-88} 就是重心（质心）坐标公式 $x_{C}=\sum P_{i}x_{i}\big/\sum P_{i}$：重心是各部分重心的加权平均，权重是各部分的重力。学生甲的逐对合成是这个公式的“逐步”执行，故重物越多越繁琐；坐标公式则一式写就。大学普通物理把式中的重力换为质量，即得质心公式；严格说来，“重心”概念以重力场近似均匀（物体线度远小于地球半径）为前提，此时重心与质心重合。}
\stda 是的，我得承认，您的方法简单得多。

\tchr 还要注意，你的解法对杆上重物的个数非常敏感（每增加一个重物，解起来就越来越繁琐）。相反，我的解法在添加重物时并不变得更复杂：每添一个新的重物，式~\eqref{eq-88} 中只是分子和分母各增加一项。

\stdb 只用力矩条件，也能求出重心的位置吗？

\tchr 能。做法是对两个不同的点分别写出力矩的平衡条件。我们就照这个办法做：对点 $A$ 和点 $C$（见\figr{fig-65}~\drkgry{(b)}）考虑力矩条件。对点 $A$，力矩条件由式~(\ref{eq-87}) 给出；对点 $C$，方程为
\begin{equation}%
N (5a-\overline{AB}) = aP_{5}+ 2aP_{4} + 3aP_{3} + 4aP_{2} + 5aP_{1}
\label{eq-89}
\end{equation}
把式~(\ref{eq-87}) 除以式~(\ref{eq-89})，得到
\begin{equation*}%
\frac{\overline{AB}}{5a - \overline{AB}} = \frac{aP_{2} + 2aP_{3} + 3aP_{4} + 4aP_{5} + 5aP_{6}
}{aP_{5}+ 2aP_{4} + 3aP_{3} + 4aP_{2} + 5aP_{1}}
\end{equation*}
由此
%\begin{equation}%
\begin{align*}
&  \overline{AB} \left( aP_{5}+ 2aP_{4} + 3aP_{3} + 4aP_{2} + 5aP_{1} + \right. \\
&\qquad  \left. aP_{2} + 2aP_{3} + 3aP_{4} + 4aP_{5} + 5aP_{6} \right) \\
& = 5a \left( aP_{2} + 2aP_{3} + 3aP_{4} + 4aP_{5} + 5aP_{6} \right)
\end{align*}
%\end{equation*}
即
\begin{equation*}%
\begin{split}
& \overline{AB} \times 5a \left(P_{1}+P_{2}+P_{3}+P_{4}+P_{5}+P_{6} \right)\\
& = 5a \left( aP_{2} + 2aP_{3} + 3aP_{4} + 4aP_{5} + 5aP_{6} \right)
\end{split}
\end{equation*}
于是得到与式~\eqref{eq-88} 相同的结果。
\pnote{对平面平行力系，“对连线不与各力平行的两点分别取矩为零”确实可以完全代替力与力矩两个平衡条件：若合力不为零，其作用线必须同时通过两个矩心，这对竖直方向的力是不可能的（$A$、$C$ 连线沿水平方向），故合力必为零。理论力学中称这类方程组为二矩式。两式相除恰好消去 $N$——重心位置本就与那个假想的支承无关。}

\end{dialogue}

\section*{习题}
\label{prob-15}
\addcontentsline{toc}{section}{习题}


\begin{enumerate}[resume*=problems]
\item 一个圆盘上有两个圆孔，如\figr{fig-66} 所示。两孔的半径分别等于圆盘半径的二分之一和四分之一。求这个圆盘的重心。
\begin{figure}[!htb]%[-3cm]  %
\centering
\includegraphics[width=0.4\textwidth]{figs/ch-15/fig-066a.pdf}
\caption[求所给物体的重心。]{问题是求所给物体的重心。\label{fig-66}}
\end{figure}


\end{enumerate}


\cleardoublepage

\thispagestyle{empty}
\vspace*{-1.5cm}
\begin{figure*}[!ht]
\centering
\includegraphics[width=0.75\textwidth]{figs/sec/tarasov-qasp-sec-08.pdf}
\end{figure*}
%\begin{centering}
\begin{fullwidth}
%%\paragraph{}
{\textsf{\Large \hlred{阿基米德原理通常不会引起特别的注意。这是备考物理的学生常犯的一个错误。在这一原理的基础上，可以设计出非常有趣的问题和习题。我们将讨论阿基米德原理对处于失重状态的物体是否适用的问题。
}}}
\end{fullwidth}

% NEWTERMS: centre of gravity | 重心 | glossary 已有（第 2 章首现），本章沿用、不重复括注
% NEWTERMS: bar | 杆 | 本章核心场景词；第 14 章英文亦用 bar（并行翻译），如彼处译“细杆/直杆”，建议全书统一为“杆”
% NEWTERMS: load / system of loads | 重物 / 重物系统 | 重物承第 3 章定名
% NEWTERMS: rule of force moments | 力矩定则 | 即杠杆平衡条件，见本章评注；对应第 14 章 moment condition
% NEWTERMS: first / second condition of equilibrium | 第一个平衡条件 / 第二个平衡条件（力矩条件） | 与第 14 章（并行翻译）对应术语，如有出入以先定稿者为准
% NEWTERMS: moment condition | 力矩条件 | 同上，第 14 章“Second condition (moment condition)”
% NEWTERMS: negative weight | 负重量 | 首现括注英文；相应方法名“负面积（负质量）法”见评注
% NEWTERMS: circular disk / circular hole | 圆盘 / 圆孔 | 挖孔类重心题基本词
% NEWTERMS: point of application of the resultant | 合力的作用点 | 合力见 glossary
% NEWTERMS: bearing reaction | 支承力 | glossary 已有（第 2 章定名，符号 $N$），本章沿用
% NEWTERMS: Archimedes' principle | 阿基米德原理 | 章尾格言首现；正文将于下一章起展开
% NEWTERMS: weightlessness | 失重 | 第 7 章格言已首现，本章章尾格言再用
